\documentclass[]{article}
\usepackage{lmodern}
\usepackage{amssymb,amsmath}
\usepackage{ifxetex,ifluatex}
\usepackage{fixltx2e} % provides \textsubscript
\ifnum 0\ifxetex 1\fi\ifluatex 1\fi=0 % if pdftex
  \usepackage[T1]{fontenc}
  \usepackage[utf8]{inputenc}
\else % if luatex or xelatex
  \ifxetex
    \usepackage{mathspec}
  \else
    \usepackage{fontspec}
  \fi
  \defaultfontfeatures{Ligatures=TeX,Scale=MatchLowercase}
\fi
% use upquote if available, for straight quotes in verbatim environments
\IfFileExists{upquote.sty}{\usepackage{upquote}}{}
% use microtype if available
\IfFileExists{microtype.sty}{%
\usepackage[]{microtype}
\UseMicrotypeSet[protrusion]{basicmath} % disable protrusion for tt fonts
}{}
\PassOptionsToPackage{hyphens}{url} % url is loaded by hyperref
\usepackage[unicode=true]{hyperref}
\hypersetup{
            pdfborder={0 0 0},
            breaklinks=true}
\urlstyle{same}  % don't use monospace font for urls
\usepackage[margin=2.54cm,a4paper]{geometry}
\usepackage{color}
\usepackage{fancyvrb}
\newcommand{\VerbBar}{|}
\newcommand{\VERB}{\Verb[commandchars=\\\{\}]}
\DefineVerbatimEnvironment{Highlighting}{Verbatim}{commandchars=\\\{\}}
% Add ',fontsize=\small' for more characters per line
\newenvironment{Shaded}{}{}
\newcommand{\KeywordTok}[1]{\textcolor[rgb]{0.00,0.44,0.13}{\textbf{#1}}}
\newcommand{\DataTypeTok}[1]{\textcolor[rgb]{0.56,0.13,0.00}{#1}}
\newcommand{\DecValTok}[1]{\textcolor[rgb]{0.25,0.63,0.44}{#1}}
\newcommand{\BaseNTok}[1]{\textcolor[rgb]{0.25,0.63,0.44}{#1}}
\newcommand{\FloatTok}[1]{\textcolor[rgb]{0.25,0.63,0.44}{#1}}
\newcommand{\ConstantTok}[1]{\textcolor[rgb]{0.53,0.00,0.00}{#1}}
\newcommand{\CharTok}[1]{\textcolor[rgb]{0.25,0.44,0.63}{#1}}
\newcommand{\SpecialCharTok}[1]{\textcolor[rgb]{0.25,0.44,0.63}{#1}}
\newcommand{\StringTok}[1]{\textcolor[rgb]{0.25,0.44,0.63}{#1}}
\newcommand{\VerbatimStringTok}[1]{\textcolor[rgb]{0.25,0.44,0.63}{#1}}
\newcommand{\SpecialStringTok}[1]{\textcolor[rgb]{0.73,0.40,0.53}{#1}}
\newcommand{\ImportTok}[1]{#1}
\newcommand{\CommentTok}[1]{\textcolor[rgb]{0.38,0.63,0.69}{\textit{#1}}}
\newcommand{\DocumentationTok}[1]{\textcolor[rgb]{0.73,0.13,0.13}{\textit{#1}}}
\newcommand{\AnnotationTok}[1]{\textcolor[rgb]{0.38,0.63,0.69}{\textbf{\textit{#1}}}}
\newcommand{\CommentVarTok}[1]{\textcolor[rgb]{0.38,0.63,0.69}{\textbf{\textit{#1}}}}
\newcommand{\OtherTok}[1]{\textcolor[rgb]{0.00,0.44,0.13}{#1}}
\newcommand{\FunctionTok}[1]{\textcolor[rgb]{0.02,0.16,0.49}{#1}}
\newcommand{\VariableTok}[1]{\textcolor[rgb]{0.10,0.09,0.49}{#1}}
\newcommand{\ControlFlowTok}[1]{\textcolor[rgb]{0.00,0.44,0.13}{\textbf{#1}}}
\newcommand{\OperatorTok}[1]{\textcolor[rgb]{0.40,0.40,0.40}{#1}}
\newcommand{\BuiltInTok}[1]{#1}
\newcommand{\ExtensionTok}[1]{#1}
\newcommand{\PreprocessorTok}[1]{\textcolor[rgb]{0.74,0.48,0.00}{#1}}
\newcommand{\AttributeTok}[1]{\textcolor[rgb]{0.49,0.56,0.16}{#1}}
\newcommand{\RegionMarkerTok}[1]{#1}
\newcommand{\InformationTok}[1]{\textcolor[rgb]{0.38,0.63,0.69}{\textbf{\textit{#1}}}}
\newcommand{\WarningTok}[1]{\textcolor[rgb]{0.38,0.63,0.69}{\textbf{\textit{#1}}}}
\newcommand{\AlertTok}[1]{\textcolor[rgb]{1.00,0.00,0.00}{\textbf{#1}}}
\newcommand{\ErrorTok}[1]{\textcolor[rgb]{1.00,0.00,0.00}{\textbf{#1}}}
\newcommand{\NormalTok}[1]{#1}
\IfFileExists{parskip.sty}{%
\usepackage{parskip}
}{% else
\setlength{\parindent}{0pt}
\setlength{\parskip}{6pt plus 2pt minus 1pt}
}
\setlength{\emergencystretch}{3em}  % prevent overfull lines
\providecommand{\tightlist}{%
  \setlength{\itemsep}{0pt}\setlength{\parskip}{0pt}}
\setcounter{secnumdepth}{0}
% Redefines (sub)paragraphs to behave more like sections
\ifx\paragraph\undefined\else
\let\oldparagraph\paragraph
\renewcommand{\paragraph}[1]{\oldparagraph{#1}\mbox{}}
\fi
\ifx\subparagraph\undefined\else
\let\oldsubparagraph\subparagraph
\renewcommand{\subparagraph}[1]{\oldsubparagraph{#1}\mbox{}}
\fi

% set default figure placement to htbp
\makeatletter
\def\fps@figure{htbp}
\makeatother

% Pandoc PDF style for Chinese survey reports
\usepackage{fontspec}
\usepackage{xeCJK}
\usepackage{indentfirst}
\usepackage{titlesec}
\usepackage{graphicx}
\usepackage{caption}
\usepackage{float}
\usepackage{etoolbox}

% ==================== 字体配置 ====================
% 设置中文字体 - 使用跨平台兼容的字体配置
% 优先使用系统自带字体，如果不存在则使用备用字体

% 黑体（标题用）
\IfFontExistsTF{SimHei}{
  \setCJKmainfont{SimHei}
  \newCJKfontfamily\fontheiti{SimHei}
}{
  \IfFontExistsTF{Heiti SC}{
    \setCJKmainfont{Heiti SC}
    \newCJKfontfamily\fontheiti{Heiti SC}
  }{
    \IfFontExistsTF{Heiti TC}{
      \setCJKmainfont{Heiti TC}
      \newCJKfontfamily\fontheiti{Heiti TC}
    }{
      \IfFontExistsTF{WenQuanYi Micro Hei}{
        \setCJKmainfont{WenQuanYi Micro Hei}
        \newCJKfontfamily\fontheiti{WenQuanYi Micro Hei}
      }{
        \IfFontExistsTF{FandolHei-Regular}{
          \setCJKmainfont{FandolHei-Regular}
          \newCJKfontfamily\fontheiti{FandolHei-Regular}
        }{
          \setCJKmainfont{Arial Unicode MS}
          \newCJKfontfamily\fontheiti{Arial Unicode MS}
        }
      }
    }
  }
}

% 宋体（正文用）
\IfFontExistsTF{SimSun}{
  \newCJKfontfamily\fontsongti{SimSun}
}{
  \IfFontExistsTF{Songti SC}{
    \newCJKfontfamily\fontsongti{Songti SC}
  }{
    \IfFontExistsTF{Songti TC}{
      \newCJKfontfamily\fontsongti{Songti TC}
    }{
      \IfFontExistsTF{WenQuanYi Micro Hei Mono}{
        \newCJKfontfamily\fontsongti{WenQuanYi Micro Hei Mono}
      }{
        \IfFontExistsTF{FandolSong-Regular}{
          \newCJKfontfamily\fontsongti{FandolSong-Regular}
        }{
          \newCJKfontfamily\fontsongti{Arial Unicode MS}
        }
      }
    }
  }
}

% 楷体（强调用）
\IfFontExistsTF{Kaiti}{
  \newCJKfontfamily\fontkaiti{Kaiti}
}{
  \IfFontExistsTF{Kaiti SC}{
    \newCJKfontfamily\fontkaiti{Kaiti SC}
  }{
    \IfFontExistsTF{Kaiti TC}{
      \newCJKfontfamily\fontkaiti{Kaiti TC}
    }{
      \IfFontExistsTF{FandolKai-Regular}{
        \newCJKfontfamily\fontkaiti{FandolKai-Regular}
      }{
        \IfFontExistsTF{WenQuanYi Micro Hei}{
          \newCJKfontfamily\fontkaiti{WenQuanYi Micro Hei}
        }{
          \newCJKfontfamily\fontkaiti{SimSun}
        }
      }
    }
  }
}

% 仿宋（公文用）
\IfFontExistsTF{FangSong}{
  \newCJKfontfamily\fontfangsong{FangSong}
}{
  \IfFontExistsTF{STFangsong}{
    \newCJKfontfamily\fontfangsong{STFangsong}
  }{
    \IfFontExistsTF{FandolFang-Regular}{
      \newCJKfontfamily\fontfangsong{FandolFang-Regular}
    }{
      \IfFontExistsTF{WenQuanYi Micro Hei Mono}{
        \newCJKfontfamily\fontfangsong{WenQuanYi Micro Hei Mono}
      }{
        \newCJKfontfamily\fontfangsong{SimSun}
      }
    }
  }
}

% 设置英文字体 - 增加更多备选字体
\IfFontExistsTF{Times New Roman}{
  \setmainfont{Times New Roman}
}{
  \IfFontExistsTF{Times}{
    \setmainfont{Times}
  }{
    \setmainfont{DejaVu Serif}
  }
}

\IfFontExistsTF{Arial}{
  \setsansfont{Arial}
}{
  \IfFontExistsTF{Helvetica}{
    \setsansfont{Helvetica}
  }{
    \setsansfont{DejaVu Sans}
  }
}

\IfFontExistsTF{Courier New}{
  \setmonofont{Courier New}
}{
  \IfFontExistsTF{Courier}{
    \setmonofont{Courier}
  }{
    \setmonofont{DejaVu Sans Mono}
  }
}

% ==================== 页面设置 ====================
% 正文：仿宋，三号（16pt），固定行距 28pt
\AtBeginDocument{\fontfangsong\fontsize{16pt}{28pt}\selectfont}
\setlength{\parindent}{2em}
\setlength{\parskip}{0pt}

% 不显示默认标题
\AtBeginDocument{\let\maketitle\relax}

% ==================== 表格设置 ====================
% 表格：宋体，四号（14pt），单倍行距
\AtBeginEnvironment{tabular}{\fontsongti\fontsize{14pt}{14pt}\selectfont\renewcommand{\arraystretch}{1}}
\AtBeginEnvironment{longtable}{\fontsongti\fontsize{14pt}{14pt}\selectfont\renewcommand{\arraystretch}{1}}

% ==================== 标题样式 ====================
% #      -> 总标题：黑体，二号（22pt），居中
% ##     -> 一级标题：黑体，三号（16pt）
% ###    -> 二级标题：楷体，三号（16pt）
% ####   -> 三级标题：仿宋加粗，三号（16pt）

\titleformat{\section}
  {\centering\fontheiti\fontsize{22pt}{28pt}\selectfont}
  {\thesection}{0.8em}{}
\titleformat{\subsection}
  {\fontheiti\fontsize{20pt}{28pt}\selectfont}
  {\thesubsection}{0.6em}{}
\titleformat{\subsubsection}
  {\fontkaiti\fontsize{16pt}{28pt}\selectfont}
  {\thesubsubsection}{0.6em}{}
\titleformat{\paragraph}[hang]
  {\fontfangsong\bfseries\fontsize{16pt}{28pt}\selectfont}
  {\theparagraph}{0.6em}{}

\titlespacing*{\section}{0pt}{12pt}{12pt}
\titlespacing*{\subsection}{0pt}{10pt}{8pt}
\titlespacing*{\subsubsection}{0pt}{8pt}{6pt}
\titlespacing*{\paragraph}{0pt}{6pt}{4pt}

% ==================== 页眉页脚 ====================
% 仅显示页脚页码，无页眉线
\pagestyle{plain}
\raggedbottom
\makeatletter
\let\ps@plain\ps@empty
\makeatother
\usepackage{fancyhdr}
\fancypagestyle{plain}{
  \fancyhf{}
  \fancyfoot[C]{\thepage}
  \renewcommand{\headrulewidth}{0pt}
  \renewcommand{\footrulewidth}{0pt}
}

% ==================== 图片设置 ====================
\floatplacement{figure}{htbp}
\renewcommand{\topfraction}{0.9}
\renewcommand{\bottomfraction}{0.8}
\renewcommand{\textfraction}{0.07}
\renewcommand{\floatpagefraction}{0.8}
\setlength{\textfloatsep}{10pt plus 2pt minus 2pt}
\setlength{\intextsep}{10pt plus 2pt minus 2pt}
\setkeys{Gin}{width=0.9\linewidth,keepaspectratio}
\captionsetup[figure]{font={small},justification=centering,singlelinecheck=false,skip=6pt}

\date{}

\begin{document}

\section{Python AI
智能体开发学习项目}\label{python-ai-ux667aux80fdux4f53ux5f00ux53d1ux5b66ux4e60ux9879ux76ee}

\subsection{项目简介}\label{ux9879ux76eeux7b80ux4ecb}

这是一个基于 Python 的 AI 智能体开发学习项目，旨在帮助学习者了解如何使用
Python 标准库与 OpenAI 兼容协议的 LLM（大语言模型）进行交互。

\subsection{目录结构}\label{ux76eeux5f55ux7ed3ux6784}

\begin{verbatim}
trae111/
├── venv/                    # 虚拟环境
├── practice01/              # 练习1：LLM 基础访问
│   ├── __init__.py
│   └── llm_client.py
├── practice02/              # 练习2：工具调用
│   ├── __init__.py
│   ├── file_tools.py
│   └── llm_toolcall.py
├── practice03/              # 练习3：流式输出聊天
│   ├── __init__.py
│   └── llm_chat.py
├── practice04/              # 练习4：聊天记录总结 + 关键信息提取 + 聊天历史搜索
│   ├── __init__.py
│   ├── llm_chat_summary.py
│   └── llm_chat_analyzer.py
├── practice05/              # 练习5：聊天记录总结 + 关键信息提取 + 聊天历史搜索 + AnythingLLM 查询
│   ├── __init__.py
│   └── llm_chat_analyzer.py
├── .env                     # 环境变量文件（用户创建）
├── .env.example             # 环境变量模板
├── .gitignore               # Git 忽略文件
├── requirements.txt         # 项目依赖
└── README.md                # 项目说明文档
\end{verbatim}

\subsection{环境配置}\label{ux73afux5883ux914dux7f6e}

\begin{enumerate}
\def\labelenumi{\arabic{enumi}.}
\item
  \textbf{虚拟环境已创建}： \texttt{bash\ \ \ \ python\ -m\ venv\ venv}
\item
  \textbf{激活虚拟环境}：
\end{enumerate}

\begin{itemize}
\tightlist
\item
  Windows PowerShell:
  \texttt{.\textbackslash{}venv\textbackslash{}Scripts\textbackslash{}Activate.ps1}
\item
  Windows CMD:
  \texttt{venv\textbackslash{}Scripts\textbackslash{}activate.bat}
\item
  Linux/Mac: \texttt{source\ venv/bin/activate}
\end{itemize}

\begin{enumerate}
\def\labelenumi{\arabic{enumi}.}
\setcounter{enumi}{2}
\item
  \textbf{安装依赖}：
  \texttt{bash\ \ \ \ pip\ install\ -r\ requirements.txt}
\item
  \textbf{配置环境变量}：
\end{enumerate}

\begin{itemize}
\tightlist
\item
  复制 \texttt{env.example} 为 \texttt{.env} 文件
\item
  填写正确的 LLM 配置信息：
  \texttt{LLM\_BASE\_URL=http://127.0.0.1:1234/v1\ \ LLM\_MODEL=qwen3.5-4b\ \ LLM\_API\_KEY=your\_api\_key\_here\ \ LLM\_TIMEOUT=3600}
\end{itemize}

\subsection{练习模块说明}\label{ux7ec3ux4e60ux6a21ux5757ux8bf4ux660e}

\subsubsection{Practice01: LLM
基础客户端}\label{practice01-llm-ux57faux7840ux5ba2ux6237ux7aef}

\textbf{文件位置}: \texttt{practice01/llm\_client.py}

\textbf{功能说明}: - 读取项目根目录的 \texttt{.env} 文件加载配置 - 使用
Python 标准 HTTP 库（urllib）访问 LLM API - 支持 OpenAI
兼容协议的聊天接口 - 统计 token 消耗、响应时间和处理速度

\textbf{运行方法}:

\begin{Shaded}
\begin{Highlighting}[]
\ExtensionTok{python}\NormalTok{ practice01/llm_client.py}
\end{Highlighting}
\end{Shaded}

\textbf{核心功能}: 1. \textbf{环境变量加载} (\texttt{load\_env\_file}):
解析 \texttt{.env} 文件，提取 LLM 配置 2. \textbf{API 调用}
(\texttt{call\_llm\_api}): 发送 HTTP POST 请求到 LLM 服务 3.
\textbf{性能统计}: 自动计算 token 数量、响应时间、tokens/秒 4.
\textbf{错误处理}: 捕获 HTTP 错误、URL 错误等异常情况

\textbf{输出示例}:

\begin{verbatim}
============================================================
✅ LLM 响应成功

响应内容:
------------------------------------------------------------
你好！我是一个 AI 助手，很高兴为你服务。

------------------------------------------------------------
📊 性能统计:
  总 Token 数: 25
  输入 Token: 15
  输出 Token: 10
  响应时间: 2.35 秒
  Token 处理速度: 4.26 tokens/秒
============================================================
\end{verbatim}

\textbf{教学目标}: 1. 学习如何使用 Python 标准库进行 HTTP 请求 2. 了解
OpenAI 兼容 API 的请求/响应格式 3. 掌握环境变量的管理和使用 4.
学习基本的性能统计和错误处理

\subsubsection{Practice02:
工具调用客户端}\label{practice02-ux5de5ux5177ux8c03ux7528ux5ba2ux6237ux7aef}

\textbf{文件位置}: \texttt{practice02/}

\textbf{功能说明}: - 实现 6 个工具：5 个文件操作工具 + 1 个网络访问工具
- 将工具定义作为系统提示词发送给 LLM - LLM
可以根据用户需求调用相应的工具 - 支持工具调用解析和执行

\textbf{运行方法}:

\begin{Shaded}
\begin{Highlighting}[]
\ExtensionTok{python}\NormalTok{ practice02/llm_toolcall.py}
\end{Highlighting}
\end{Shaded}

\textbf{核心功能}: 1. \textbf{文件操作工具} (\texttt{file\_tools.py}): -
\texttt{list\_files}: 列出目录下的所有文件（包括属性和大小） -
\texttt{rename\_file}: 重命名文件 - \texttt{delete\_file}: 删除文件 -
\texttt{create\_file}: 创建新文件并写入内容 - \texttt{read\_file}:
读取文件内容

\begin{enumerate}
\def\labelenumi{\arabic{enumi}.}
\setcounter{enumi}{1}
\tightlist
\item
  \textbf{网络访问工具} (\texttt{file\_tools.py}):
\end{enumerate}

\begin{itemize}
\tightlist
\item
  \texttt{curl\_web}: 通过 curl 访问网页并返回网页内容
\item
  支持 HTTP 和 HTTPS 协议
\item
  自动处理编码和错误
\end{itemize}

\begin{enumerate}
\def\labelenumi{\arabic{enumi}.}
\setcounter{enumi}{2}
\tightlist
\item
  \textbf{工具调用} (\texttt{llm\_toolcall.py}):
\end{enumerate}

\begin{itemize}
\tightlist
\item
  \texttt{parse\_tool\_call}: 解析 LLM 返回的工具调用请求
\item
  \texttt{execute\_tool\_call}: 执行工具并返回结果
\item
  \texttt{get\_tool\_definition}: 生成工具定义的 JSON 格式
\end{itemize}

\begin{enumerate}
\def\labelenumi{\arabic{enumi}.}
\setcounter{enumi}{3}
\tightlist
\item
  \textbf{交互式对话}: 支持多轮对话，自动维护聊天历史
\end{enumerate}

\textbf{使用示例}:

\begin{verbatim}
> 列出当前目录的文件
🔧 检测到工具调用: list_files
   参数: {"directory": "."}

============================================================
🔧 工具执行结果
============================================================
✅ 工具执行成功
   找到 5 个文件:
     - __init__.py (file, 0.00 B)
     - file_tools.py (file, 12.34 KB)
     - llm_toolcall.py (file, 15.67 KB)

> 访问 https://www.example.com
🔧 检测到工具调用: curl_web
   参数: {"url": "https://www.example.com"}

============================================================
🔧 工具执行结果
============================================================
✅ 工具执行成功
   URL: https://www.example.com
   内容长度: 1256 字节
   状态码: 200
\end{verbatim}

\textbf{教学目标}: 1. 学习如何设计和实现工具调用机制 2.
掌握工具定义和参数验证 3. 了解如何让 LLM 理解和使用工具 4.
学习工具调用的解析和执行流程 5. 掌握网络请求和网页内容获取技术

\subsubsection{Practice03:
流式输出聊天客户端}\label{practice03-ux6d41ux5f0fux8f93ux51faux804aux5929ux5ba2ux6237ux7aef}

\textbf{文件位置}: \texttt{practice03/llm\_chat.py}

\textbf{功能说明}: - 支持终端界面输入聊天内容 - 实现流式输出，实时显示
AI 响应 - 自动维护聊天历史，添加到上下文 -
支持特殊命令：clear（清空历史）、history（查看历史） - 按 Ctrl+C
退出程序

\textbf{运行方法}:

\begin{Shaded}
\begin{Highlighting}[]
\ExtensionTok{python}\NormalTok{ practice03/llm_chat.py}
\end{Highlighting}
\end{Shaded}

\textbf{核心功能}: 1. \textbf{流式输出}
(\texttt{call\_llm\_api\_stream}): - 使用 SSE (Server-Sent Events)
协议接收流式响应 - 实时输出 AI 生成的文本 - 提升用户体验，减少等待感

\begin{enumerate}
\def\labelenumi{\arabic{enumi}.}
\setcounter{enumi}{1}
\tightlist
\item
  \textbf{聊天历史管理}:
\end{enumerate}

\begin{itemize}
\tightlist
\item
  自动保存用户和 AI 的对话
\item
  每次请求都包含完整的历史记录
\item
  支持清空历史和查看历史
\end{itemize}

\begin{enumerate}
\def\labelenumi{\arabic{enumi}.}
\setcounter{enumi}{2}
\tightlist
\item
  \textbf{性能统计}:
\end{enumerate}

\begin{itemize}
\tightlist
\item
  实时显示 token 消耗
\item
  统计响应时间和处理速度
\item
  显示当前历史记录消息数
\end{itemize}

\textbf{使用示例}:

\begin{verbatim}
============================================================
🚀 Practice03: 流式输出聊天客户端
============================================================

💡 提示:
   - 输入消息与 AI 对话
   - 输入 'clear' 清空聊天历史
   - 输入 'history' 查看聊天历史
   - 按 Ctrl+C 退出程序
============================================================

💬 开始聊天（按 Ctrl+C 退出）:
============================================================

你: 你好
AI: 你好！很高兴见到你。我是一个 AI 助手，有什么我可以帮助你的吗？

------------------------------------------------------------
📊 总 Token 数: 32
   响应时间: 1.23 秒
   Token 处理速度: 18.52 tokens/秒
   当前历史记录消息数: 2
------------------------------------------------------------
\end{verbatim}

\textbf{教学目标}: 1. 学习如何实现流式输出 2. 掌握 SSE 协议的使用 3.
了解如何维护聊天上下文 4. 学习终端交互程序的设计

\subsubsection{Practice04:
聊天记录总结}\label{practice04-ux804aux5929ux8bb0ux5f55ux603bux7ed3}

\textbf{文件位置}: \texttt{practice04/llm\_chat\_summary.py}

\textbf{功能说明}: -
自动检测聊天历史，当超过5轮对话或上下文长度超过3k字符时触发总结 -
对前70\%左右的对话内容进行压缩总结 -
最后30\%的对话保留原文，确保最新信息不丢失 -
将总结内容作为系统提示词的一部分，保持对话连贯性

\textbf{运行方法}:

\begin{Shaded}
\begin{Highlighting}[]
\ExtensionTok{python}\NormalTok{ practice04/llm_chat_summary.py}
\end{Highlighting}
\end{Shaded}

\textbf{核心功能}: 1. \textbf{自动检测机制}
(\texttt{should\_summarize}): - 检测对话轮数（超过5轮触发） -
检测上下文字符长度（超过3000字符触发） - 智能判断是否需要执行总结

\begin{enumerate}
\def\labelenumi{\arabic{enumi}.}
\setcounter{enumi}{1}
\tightlist
\item
  \textbf{聊天总结} (\texttt{summarize\_chat\_history}):
\end{enumerate}

\begin{itemize}
\tightlist
\item
  将前70\%的对话发送给 LLM 生成摘要
\item
  保留后30\%的对话原文
\item
  生成包含主题、关键信息、用户需求、重要结论的摘要
\end{itemize}

\begin{enumerate}
\def\labelenumi{\arabic{enumi}.}
\setcounter{enumi}{2}
\tightlist
\item
  \textbf{上下文压缩}:
\end{enumerate}

\begin{itemize}
\tightlist
\item
  原始消息 → 压缩后消息
\item
  保留系统消息和最新对话
\item
  显著减少 token 消耗
\end{itemize}

\textbf{使用示例}:

\begin{verbatim}
============================================================
🚀 Practice04: 聊天记录总结功能
============================================================

💡 提示:
   - 当聊天超过5轮或上下文长度超过3k时，自动执行总结
   - 前70%左右的内容会被压缩，最后30%保留原文
   - 输入 'clear' 清空聊天历史
   - 输入 'history' 查看聊天历史
   - 按 Ctrl+C 退出程序
============================================================

💬 开始聊天（按 Ctrl+C 退出）:
============================================================

你: 你好
AI: 你好！很高兴见到你。

...

你: 继续讨论AI的发展
AI: 当然可以，AI的发展确实非常迅速...

============================================================
💡 对话轮数 6 超过限制 5
============================================================

📝 检测到聊天历史过长，正在执行总结压缩...
============================================================

📊 总结统计:
   总消息数: 12
   系统消息: 1
   需要压缩: 8
   保留原文: 3

🤖 正在调用 LLM 生成摘要...

📋 生成的摘要:
------------------------------------------------------------
1. 对话主题：AI助手的能力和AI技术的发展
2. 关键信息：用户了解了AI助手的功能和AI技术的发展历程
3. 用户需求：讨论AI相关话题
4. 重要结论：AI技术发展迅速，未来潜力巨大
------------------------------------------------------------

✅ 总结完成！
   原始消息数: 12
   压缩后消息数: 5
   压缩率: 58.3%

✅ 已将 12 条消息压缩为 5 条

📊 总 Token 数: 150
   响应时间: 5.23 秒
   Token 处理速度: 18.52 tokens/秒
   当前对话轮数: 3
   当前上下文长度: 1250
------------------------------------------------------------
\end{verbatim}

\textbf{教学目标}: 1. 学习如何实现聊天历史的自动管理和压缩 2.
掌握对话上下文的智能控制技术 3. 了解 LLM 在总结任务中的应用 4.
学习对话轮数和上下文字符长度的计算方法 5. 掌握聊天记录的结构化压缩策略

\subsubsection{Practice04: 聊天记录总结 + 关键信息提取 +
聊天历史搜索}\label{practice04-ux804aux5929ux8bb0ux5f55ux603bux7ed3-ux5173ux952eux4fe1ux606fux63d0ux53d6-ux804aux5929ux5386ux53f2ux641cux7d22}

\textbf{文件位置}: \texttt{practice04/llm\_chat\_analyzer.py}

\textbf{功能说明}: -
自动检测聊天历史，当超过5轮对话或上下文长度超过3k字符时触发总结 -
对前70\%左右的对话内容进行压缩总结，最后30\%的对话保留原文 -
每5次聊天自动提取关键信息，按照5W规则（Who, What, When, Where,
Why）进行分析 - 将关键信息增量保存到
\texttt{D:\textbackslash{}chat-log\textbackslash{}log.txt} 文件 - 支持
\texttt{/search} 命令搜索聊天历史，LLM 会根据日志文件和用户查询生成回答

\textbf{运行方法}:

\begin{Shaded}
\begin{Highlighting}[]
\ExtensionTok{python}\NormalTok{ practice04/llm_chat_analyzer.py}
\end{Highlighting}
\end{Shaded}

\textbf{核心功能}: 1. \textbf{自动检测机制}
(\texttt{should\_summarize}): - 检测对话轮数（超过5轮触发） -
检测上下文字符长度（超过3000字符触发） - 智能判断是否需要执行总结

\begin{enumerate}
\def\labelenumi{\arabic{enumi}.}
\setcounter{enumi}{1}
\tightlist
\item
  \textbf{聊天总结} (\texttt{summarize\_chat\_history}):
\end{enumerate}

\begin{itemize}
\tightlist
\item
  将前70\%的对话发送给 LLM 生成摘要
\item
  保留后30\%的对话原文
\item
  生成包含主题、关键信息、用户需求、重要结论的摘要
\end{itemize}

\begin{enumerate}
\def\labelenumi{\arabic{enumi}.}
\setcounter{enumi}{2}
\tightlist
\item
  \textbf{关键信息提取} (\texttt{extract\_key\_info}):
\end{enumerate}

\begin{itemize}
\tightlist
\item
  每5次聊天自动触发
\item
  按照5W规则提取：Who（谁）、What（做了什么事）、When（什么时候，可选）、Where（在何处，可选）、Why（为什么要做这个事，可选）
\item
  支持提取多条关键信息
\end{itemize}

\begin{enumerate}
\def\labelenumi{\arabic{enumi}.}
\setcounter{enumi}{3}
\tightlist
\item
  \textbf{日志管理} (\texttt{save\_key\_info\_to\_log}):
\end{enumerate}

\begin{itemize}
\tightlist
\item
  自动创建 \texttt{D:\textbackslash{}chat-log} 目录
\item
  增量更新 \texttt{log.txt} 文件
\item
  每条记录包含时间戳和关键信息
\end{itemize}

\begin{enumerate}
\def\labelenumi{\arabic{enumi}.}
\setcounter{enumi}{4}
\tightlist
\item
  \textbf{聊天历史搜索} (\texttt{search\_chat\_history}):
\end{enumerate}

\begin{itemize}
\tightlist
\item
  支持 \texttt{/search} 命令
\item
  读取 \texttt{log.txt} 文件
\item
  LLM 根据日志内容和用户查询生成回答
\end{itemize}

\textbf{使用示例}:

\begin{verbatim}
============================================================
🚀 Practice04: 聊天记录总结 + 关键信息提取 + 聊天历史搜索
============================================================

💡 提示:
   - 当聊天超过5轮或上下文长度超过3k时，自动执行总结
   - 每5次聊天自动提取关键信息（5W规则）并保存到日志
   - 输入 '/search 查询内容' 搜索聊天历史
   - 输入 'clear' 清空聊天历史
   - 输入 'history' 查看聊天历史
   - 按 Ctrl+C 退出程序
============================================================

💬 开始聊天（按 Ctrl+C 退出）:
============================================================

你: 你好，我是小明
AI: 你好小明！很高兴认识你。有什么我可以帮助你的吗？

你: 我想学习Python编程
AI: 很好的选择！Python是一门非常流行的编程语言...

...

你: 继续讨论Python
AI: 当然可以，Python有很多应用领域...

============================================================
💡 已完成 5 次聊天，正在提取关键信息...
============================================================

🔍 正在提取关键信息（5W规则）...
============================================================

🤖 正在调用 LLM 提取关键信息...

📋 提取的关键信息:
------------------------------------------------------------
【关键信息 1】
- Who（谁）：小明
- What（做了什么事）：想要学习Python编程
- When（什么时候）：未明确
- Where（在何处）：未明确
- Why（为什么要做这个事）：想要学习编程技能
------------------------------------------------------------

✅ 已创建目录: D:\chat-log
✅ 关键信息已保存到: D:\chat-log\log.txt
✅ 关键信息已成功保存

你: /search 小明学过什么
============================================================
🔍 正在搜索聊天历史...
============================================================

📂 已读取日志文件: D:\chat-log\log.txt
📊 日志文件大小: 1234 字符

🤖 正在调用 LLM 搜索...

📋 搜索结果:
------------------------------------------------------------
根据聊天历史记录，小明想要学习Python编程。在对话中，小明表达了学习Python的兴趣，AI向他介绍了Python的特点和应用领域。
------------------------------------------------------------
\end{verbatim}

\textbf{教学目标}: 1. 学习如何实现聊天历史的自动管理和压缩 2.
掌握对话上下文的智能控制技术 3. 了解 LLM 在总结任务中的应用 4.
学习对话轮数和上下文字符长度的计算方法 5. 掌握聊天记录的结构化压缩策略
6. 学习5W规则的信息提取方法 7. 掌握日志文件的增量更新技术 8.
了解如何实现聊天历史搜索功能 9. 学习如何让 LLM 理解和执行搜索任务

\subsubsection{Practice05: 聊天记录总结 + 关键信息提取 + 聊天历史搜索 +
AnythingLLM
查询}\label{practice05-ux804aux5929ux8bb0ux5f55ux603bux7ed3-ux5173ux952eux4fe1ux606fux63d0ux53d6-ux804aux5929ux5386ux53f2ux641cux7d22-anythingllm-ux67e5ux8be2}

\textbf{文件位置}: \texttt{practice05/llm\_chat\_analyzer.py}

\textbf{功能说明}: -
自动检测聊天历史，当超过5轮对话或上下文长度超过3k字符时触发总结 -
对前70\%左右的对话内容进行压缩总结，最后30\%的对话保留原文 -
每5次聊天自动提取关键信息，按照5W规则（Who, What, When, Where,
Why）进行分析 - 将关键信息增量保存到
\texttt{D:\textbackslash{}chat-log\textbackslash{}log.txt} 文件 - 支持
\texttt{/search} 命令搜索聊天历史 - 支持 AnythingLLM
文档仓库查询，当用户提到``文档仓库''、``文件仓库''、``仓库''等关键词时自动触发

\textbf{运行方法}:

\begin{Shaded}
\begin{Highlighting}[]
\ExtensionTok{python}\NormalTok{ practice05/llm_chat_analyzer.py}
\end{Highlighting}
\end{Shaded}

\textbf{核心功能}: 1. \textbf{自动检测机制}
(\texttt{should\_summarize}): - 检测对话轮数（超过5轮触发） -
检测上下文字符长度（超过3000字符触发） - 智能判断是否需要执行总结

\begin{enumerate}
\def\labelenumi{\arabic{enumi}.}
\setcounter{enumi}{1}
\tightlist
\item
  \textbf{聊天总结} (\texttt{summarize\_chat\_history}):
\end{enumerate}

\begin{itemize}
\tightlist
\item
  将前70\%的对话发送给 LLM 生成摘要
\item
  保留后30\%的对话原文
\item
  生成包含主题、关键信息、用户需求、重要结论的摘要
\end{itemize}

\begin{enumerate}
\def\labelenumi{\arabic{enumi}.}
\setcounter{enumi}{2}
\tightlist
\item
  \textbf{关键信息提取} (\texttt{extract\_key\_info}):
\end{enumerate}

\begin{itemize}
\tightlist
\item
  每5次聊天自动触发
\item
  按照5W规则提取：Who（谁）、What（做了什么事）、When（什么时候，可选）、Where（在何处，可选）、Why（为什么要做这个事，可选）
\item
  支持提取多条关键信息
\end{itemize}

\begin{enumerate}
\def\labelenumi{\arabic{enumi}.}
\setcounter{enumi}{3}
\tightlist
\item
  \textbf{日志管理} (\texttt{save\_key\_info\_to\_log}):
\end{enumerate}

\begin{itemize}
\tightlist
\item
  自动创建 \texttt{D:\textbackslash{}chat-log} 目录
\item
  增量更新 \texttt{log.txt} 文件
\item
  每条记录包含时间戳和关键信息
\end{itemize}

\begin{enumerate}
\def\labelenumi{\arabic{enumi}.}
\setcounter{enumi}{4}
\tightlist
\item
  \textbf{聊天历史搜索} (\texttt{search\_chat\_history}):
\end{enumerate}

\begin{itemize}
\tightlist
\item
  支持 \texttt{/search} 命令
\item
  读取 \texttt{log.txt} 文件
\item
  LLM 根据日志内容和用户查询生成回答
\end{itemize}

\begin{enumerate}
\def\labelenumi{\arabic{enumi}.}
\setcounter{enumi}{5}
\tightlist
\item
  \textbf{AnythingLLM 查询} (\texttt{anythingllm\_query}):
\end{enumerate}

\begin{itemize}
\tightlist
\item
  使用 subprocess 模块调用 curl 命令
\item
  访问
  \texttt{http://localhost:3001/api/v1/workspace/\{workspace\_slug\}/chat}
  API 接口
\item
  使用 API 密钥进行认证
\item
  支持中文编码处理
\item
  详细的错误处理和日志输出
\end{itemize}

\textbf{使用示例}:

\begin{verbatim}
============================================================
🚀 Practice05: 聊天记录总结 + 关键信息提取 + 聊天历史搜索 + AnythingLLM 查询
============================================================

💡 提示:
   - 当聊天超过5轮或上下文长度超过3k时，自动执行总结
   - 每5次聊天自动提取关键信息（5W规则）并保存到日志
   - 输入 '/search 查询内容' 搜索聊天历史
   - 提到 '文档仓库'、'文件仓库'、'仓库' 时查询 AnythingLLM
   - 输入 'clear' 清空聊天历史
   - 输入 'history' 查看聊天历史
   - 按 Ctrl+C 退出程序
============================================================

💬 开始聊天（按 Ctrl+C 退出）:
============================================================

你: 文档仓库里有什么内容？
============================================================
📚 正在查询 AnythingLLM 文档仓库...
============================================================

📡 正在访问: http://localhost:3001/api/v1/workspace/default/chat
💬 查询内容: 文档仓库里有什么内容？...

📊 命令执行结果:
   退出码: 0

📋 查询结果:
------------------------------------------------------------
文档仓库中包含以下内容：
1. 项目文档：详细介绍了项目的架构和使用方法
2. 技术文档：包含 API 文档和开发指南
3. 用户手册：指导用户如何使用系统
------------------------------------------------------------

🤖 文档仓库查询结果:
文档仓库中包含以下内容：
1. 项目文档：详细介绍了项目的架构和使用方法
2. 技术文档：包含 API 文档和开发指南
3. 用户手册：指导用户如何使用系统

你: /search 文档仓库
============================================================
🔍 正在搜索聊天历史...
============================================================

📂 已读取日志文件: D:\chat-log\log.txt
📊 日志文件大小: 1234 字符

🤖 正在调用 LLM 搜索...

📋 搜索结果:
------------------------------------------------------------
根据聊天历史记录，用户询问了文档仓库的内容，系统通过 AnythingLLM 查询到文档仓库中包含项目文档、技术文档和用户手册。
------------------------------------------------------------
\end{verbatim}

\textbf{教学目标}: 1. 学习如何实现聊天历史的自动管理和压缩 2.
掌握对话上下文的智能控制技术 3. 了解 LLM 在总结任务中的应用 4.
学习对话轮数和上下文字符长度的计算方法 5. 掌握聊天记录的结构化压缩策略
6. 学习5W规则的信息提取方法 7. 掌握日志文件的增量更新技术 8.
了解如何实现聊天历史搜索功能 9. 学习如何让 LLM 理解和执行搜索任务 10.
学习使用 subprocess 模块调用外部命令 11. 了解如何访问和使用 RESTful API
12. 掌握 API 认证和错误处理技术

\subsection{教学价值}\label{ux6559ux5b66ux4ef7ux503c}

\begin{enumerate}
\def\labelenumi{\arabic{enumi}.}
\tightlist
\item
  \textbf{API 调用实践}：学习如何通过 HTTP 请求与 LLM API 进行交互
\item
  \textbf{环境配置管理}：了解如何使用环境变量文件管理敏感配置
\item
  \textbf{性能分析}：学习如何统计和分析 API 调用的性能指标
\item
  \textbf{错误处理}：掌握基本的异常处理和错误提示
\item
  \textbf{模块化设计}：了解如何组织 Python 项目结构
\item
  \textbf{聊天历史管理}：了解如何维护和使用聊天上下文
\item
  \textbf{工具调用机制}：学习如何让 AI 使用外部工具扩展能力
\item
  \textbf{流式输出技术}：掌握实时响应和用户体验优化
\item
  \textbf{网络请求处理}：学习如何使用 Python 标准库进行 HTTP
  请求和网页内容获取
\item
  \textbf{信息提取技术}：掌握5W规则（Who, What, When, Where,
  Why）的信息提取方法
\item
  \textbf{日志管理}：学习如何实现日志文件的增量更新和管理
\item
  \textbf{聊天历史搜索}：了解如何实现基于 LLM 的智能搜索功能
\item
  \textbf{外部命令调用}：学习使用 subprocess 模块调用外部命令（如 curl）
\item
  \textbf{RESTful API 集成}：了解如何访问和使用 RESTful API
\item
  \textbf{API 认证}：掌握如何使用 API 密钥进行身份验证
\item
  \textbf{中文编码处理}：学习如何处理中文编码问题
\item
  \textbf{多系统集成}：了解如何将 LLM 与其他系统（如 AnythingLLM）集成
\end{enumerate}

\subsection{注意事项}\label{ux6ce8ux610fux4e8bux9879}

\begin{itemize}
\tightlist
\item
  请保护好您的 API 密钥，不要将其提交到版本控制系统
\item
  注意 API 调用的频率限制和 token 消耗
\item
  确保网络连接正常，以避免 API 调用失败
\item
  对于大型项目，考虑使用更高级的 HTTP 库如 requests
\item
  流式输出需要 LLM 服务支持 SSE 协议
\end{itemize}

\begin{center}\rule{0.5\linewidth}{\linethickness}\end{center}

\textbf{作者}：AI 智能体开发学习项目 \textbf{日期}：2026-04-16

\end{document}
